% TODO make hw stylesheet

\documentclass[10pt]{article}
\usepackage[letterpaper,top=1in,left=1in,right=1in]{geometry}
\usepackage[dvipsnames,svgnames]{xcolor}
\usepackage[shortlabels]{enumitem}
\usepackage[framemethod=TikZ]{mdframed}
\usepackage{amsmath,amssymb,amsthm}
\usepackage[colorlinks]{hyperref}
\usepackage{multicol}
\usepackage{tabularx}
\usepackage{bbm}
\usepackage{tikz-cd}

\hypersetup{
    colorlinks=true,
    linkcolor=blue,
    filecolor=magenta,
    urlcolor=magenta,
    pdftitle={MAT 535 HW}
}

\usepackage{mathtools}
\usepackage{thmtools}
\usepackage[textsize=footnotesize]{todonotes}
\usepackage{titling}
\usepackage{cancel}

\reversemarginpar
\mdfdefinestyle{mdbluebox}{roundcorner=10pt,innerbottommargin=9pt,
linecolor=blue,backgroundcolor=TealBlue!5,}
\declaretheoremstyle[headfont=\sffamily\bfseries\color{MidnightBlue},
mdframed={style=mdbluebox},]{thmbluebox}

\mdfdefinestyle{mdbloobox}{frametitlefont=\bfseries,innerbottommargin=8pt,
nobreak=true,backgroundcolor=TealBlue!5,linecolor=blue,}
\declaretheoremstyle[headfont=\bfseries\color{MidnightBlue},
mdframed={style=mdbloobox},headpunct={\\[3pt]},postheadspace=0pt,]{thmbloobox}

\mdfdefinestyle{mdredbox}{frametitlefont=\bfseries,innerbottommargin=8pt,
nobreak=true,backgroundcolor=Salmon!5,linecolor=RawSienna,}
\declaretheoremstyle[headfont=\bfseries\color{RawSienna},
mdframed={style=mdredbox},headpunct={\\[3pt]},postheadspace=0pt,]{thmredbox}

\mdfdefinestyle{mdgreenbox}{linecolor=ForestGreen,backgroundcolor=ForestGreen!5,
linewidth=2pt,rightline=false,leftline=true,topline=false,bottomline=false,}
\declaretheoremstyle[headfont=\bfseries\sffamily\color{ForestGreen!70!black},
mdframed={style=mdgreenbox},headpunct={ --- },]{thmgreenbox}

\mdfdefinestyle{mdorangebox}{linecolor=RedOrange,backgroundcolor=RedOrange!5,
linewidth=2pt,rightline=false,leftline=true,topline=false,bottomline=false,}
\declaretheoremstyle[headfont=\bfseries\sffamily\color{RedOrange!70!black},
mdframed={style=mdorangebox},headpunct={ --- },]{thmorangebox}

\mdfdefinestyle{mdblackbox}{linecolor=black,backgroundcolor=RedViolet!5!gray!5,
linewidth=3pt,rightline=false,leftline=true,topline=false,bottomline=false,}
\declaretheoremstyle[mdframed={style=mdblackbox}]{thmblackbox}

\declaretheoremstyle[spaceabove=6pt,spacebelow=6pt]{boldhead}
\declaretheoremstyle[spaceabove=6pt,spacebelow=6pt,headfont=\normalfont\itshape]{italichead}

\declaretheorem[style=boldhead,name=Problem]{problem}
\declaretheorem[style=boldhead,name=Problem,numbered=no]{problem*}
\declaretheorem[style=boldhead,name=Setup]{setup} % for config geo
\declaretheorem[style=boldhead,name=Example,sibling=problem]{example}
\declaretheorem[style=boldhead,name=$\bigstar$ Problem,sibling=problem]{niceprob}
\declaretheorem[style=boldhead,name=Theorem,sibling=problem]{theorem}
\declaretheorem[style=boldhead,name=Corollary,sibling=theorem]{corollary}
\declaretheorem[style=boldhead,name=Proposition,sibling=theorem]{prop}
\declaretheorem[style=boldhead,name=Lemma,numberwithin=problem]{lemma}
\declaretheorem[style=boldhead,name=Theorem,numbered=no]{theorem*}
\declaretheorem[style=boldhead,name=Lemma,numbered=no]{lemma*}
\declaretheorem[style=boldhead,name=Claim,numberwithin=problem]{claim}
\declaretheorem[style=boldhead,name=Claim,numbered=no]{claim*}
\declaretheorem[style=boldhead,name=Remark,numberwithin=problem]{remark}
\declaretheorem[style=boldhead,name=Remark,numbered=no]{remark*}
\declaretheorem[style=boldhead,name=Definition,sibling=theorem]{definition}
\declaretheorem[style=boldhead,name=Definition,numbered=no]{definition*}

\providecommand{\ol}{\overline}
\providecommand{\tensor}{\otimes}
\renewcommand{\hom}{\operatorname{Hom}}
\providecommand{\tor}{\operatorname{Tor}}
\providecommand{\Gal}{\operatorname{Gal}}
\providecommand{\ceil}[1]{\left\lceil #1 \right\rceil}
\providecommand{\floor}[1]{\left\lfloor #1 \right\rfloor}
\newcommand{\dd}{\,\mathrm{d}}
\providecommand{\ul}{\underline}
\providecommand{\eps}{\varepsilon}
\providecommand{\half}{\frac{1}{2}}
\providecommand{\CC}{\mathbb C}
\providecommand{\FF}{\mathbb F}
\providecommand{\II}{\mathbb I}
\providecommand{\NN}{\mathbb N}
\providecommand{\QQ}{\mathbb Q}
\providecommand{\RR}{\mathbb R}
\providecommand{\ZZ}{\mathbb Z}
\providecommand{\ind}{\mathbbm 1}
\providecommand{\dg}{^\circ}
\providecommand{\ii}{\item}
\providecommand{\setdiff}{\smallsetminus}
\providecommand{\alert}[1]{{\sffamily\textbf{\textcolor{blue}{#1}}}}
\providecommand{\inv}{^{-1}}
\providecommand{\mb}[1]{\mathbf{#1}}
\newcommand{\what}{\widehat}

\newenvironment{soln}{\begin{proof}[Solution]}{\end{proof}}

\providecommand{\pow}{\operatorname{Pow}}
\providecommand{\vocab}[1]{{\textbf{\textcolor{ForestGreen}{#1}}}}
\providecommand{\lcm}{\operatorname{lcm}}
\providecommand{\abs}[1]{\left\lvert #1\right\rvert}
\providecommand{\smabs}[1]{\lvert #1\rvert}
\providecommand{\innerprod}[1]{\langle\!\langle #1\rangle\!\rangle}
\providecommand{\norm}[1]{\left\| #1\right\|}
\providecommand{\smnorm}[1]{\| #1\|}
\providecommand{\gr}{\operatorname{gr}}

\usepackage{mathrsfs}

% place title at top right
\setlength{\droptitle}{-8ex}
\pretitle{\begin{flushright}\Large\bfseries}
\posttitle{\par\end{flushright}}
\preauthor{\begin{flushright}}
\postauthor{\end{flushright}}
\predate{\begin{flushright}}
\postdate{\end{flushright}}

\title{MAT 535 HW09}
\author{\large Aditya Pahuja}
\date{\large Due 04/23}
\begin{document}
\maketitle

\begin{problem}
    [\S14.6\#9]
    Determine the Galois group of $x^4 + 4x - 1$.
\end{problem}

\begin{soln}
    We can check that $p(x) = x^4 + 4x - 1$ is irreducible because
    it has no rational roots by the rational root theorem
    ($\pm 1$ are both clearly not roots) and
    it has no quadratic factors because if
    \begin{equation*}
        p(x) = (x^2 + ax + c)(x^2 + bx + d)
	= x^4 + (a + b)x^3 + (ab + c + d)x^2
	+ (ad + bc)x + cd
    \end{equation*}
    then $cd = 1$, so $\{c,d\} = \{1,-1\}$;
    $a+b = 0$, so $a = -b$;
    $ab + c + d = -a^2 = 0$, so $a = 0$;
    and $ad + bc = 4$ and $0$ simultaneously, contradiction.

    The derivative $p'(x) = 4x^3 + 4$ has one real root,
    meaning $p(x)$ as a function $\RR\to\RR$ has one global minimum
    at $x = -1$, where $p(-1) = -2$. Thus, $p(x)$
    has two real roots and two non-real roots.
    This means that the complex conjugation automorphism
    is a transposition of the two non-real roots.

    The resolvent cubic of $p(x)$ is
    \begin{equation*}
        q(x) = x^3 + 4x - 16 = (x - 2)(x^2 + 2x + 8).
    \end{equation*}
    This means the Galois group of $p(x)$
    is either $\ZZ/(4)$ or $D_8$
    by the classification in the text,
    but the former is impossible because cyclic subgroups of $S_4$
    with order $4$ do not have any transpositions,
    so the Galois group of $p(x)$ is dihedral with order $8$.
\end{soln}

\begin{problem}
    [\S14.6\#19]
    Let $f(x)$ be an irreducible polynomial of degree $4$
    in $\QQ[x]$ with discriminant $D$.
    Let $K$ denote the splitting field of $f(x)$,
    viewed as a subfield of $\CC$.
    \begin{enumerate}[(a)]
	\ii Prove that $\QQ[\sqrt D]\subseteq K$.
	\ii Let $\tau$ denote complex conjugation
	and let $\tau_K$ denote the restriction of complex conjugation to $K$.
	Prove that $\tau_K$ is an element of $\Gal(K/\QQ)$
	with order $1$ or $2$
	depending on whether every element of $K$ is real or not.
	\ii Prove that if $D < 0$ then $\Gal(K/\QQ)$
	cannot be cyclic of order $4$.
	\ii Prove generally that $\QQ(\sqrt D)$ for squarefree $D < 0$
	is not a subfield of a cyclic quartic field.
    \end{enumerate}
\end{problem}

\begin{soln}
    (a) Let $\alpha_1$, $\alpha_2$, $\alpha_3$, $\alpha_4$
    be the roots of $f$. By definition,
    \begin{equation*}
        \sqrt D = (\alpha_1-\alpha_2)
	(\alpha_1-\alpha_3)
	(\alpha_1-\alpha_4)
	(\alpha_2-\alpha_3)
	(\alpha_2-\alpha_4)
	(\alpha_3-\alpha_4)
    \end{equation*}
    which is a product of differences of elements of $K$,
    hence is in $K$, so $\QQ[\sqrt D]\subseteq K$.
    \\

    (b) Real numbers are fixed under complex conjugation,
    so if every element of $K$ is real
    then $\tau_K$ fixes every element of $K$, hence has order $1$.
    If $K$ has any non-real elements
    then $\tau_K$ is not of order $1$,
    but its order is at most $2$ because
    $\tau$ has order $2$ as an automorphism of $\CC$,
    so $\tau_K$ must have order $2$.
    \\

    (c) If $\Gal(K/\QQ)$ has order $4$, then
    $K/\QQ[\sqrt D]$ and $K/K_+$ are distinct degree $2$ subextensions
    of $K/\QQ$, where $K_+ = \RR\cap K$
    (they are distinct because $K_+\subseteq\RR$
    and $\QQ[\sqrt D]$ is not a subset of $\RR$),
    which means $\Gal(K/\QQ)$ has more than one element of order $2$,
    so it is not cyclic.
    \\

    (d) Again, if $K/\QQ$ has degree $4$,
    then $\QQ[\sqrt D]$ and $K_+$
    are distinct subfields that both correspond to
    order $2$ elements of $\Gal(K/\QQ)$,
    so the Galois group is not cyclic because it has too many
    elements of order $2$.
\end{soln}

\begin{problem}
    [\S14.6\#28]
    Let $\alpha$ be a root of the irreducible polynomial
    $f(x)\in F[x]$ and let $K = F[\alpha]$.
    Let $D$ be the discriminant of $f(x)$.
    Prove that $D = (-1)^{n(n-1)/2} N_{K/F}(f'(\alpha))$.
\end{problem}

\begin{soln}
    If $\alpha$ is not a separable element of $F$,
    then $\alpha$ has a multiplicity higher than $1$
    as a root of $f$, so $D = 0$ and $f'(\alpha) = 0$
    implies that $N_{K/F}(f'(\alpha))$, as desired.

    Assume now that $\alpha$, hence $f(x)$, is separable, and
    let $\alpha = \alpha_1$, \dots, $\alpha_n$ be the roots of $f(x)$
    in some splitting extension $L/K$ of $f(x)$. Then
    \begin{equation*}
	D = \prod_{i < j} (\alpha_i - \alpha_j)^2
	= (-1)^{n(n-1)/2}
	\prod_{i\ne j} (\alpha_i - \alpha_j)
	= (-1)^{n(n-1)/2}
	\prod_{i=1}^n \prod_{j\ne i} (\alpha_i - \alpha_j),
    \end{equation*}
    where the coefficient of $(-1)^{n(n-1)/2}$
    comes from negating $1 + 2 + \cdots + n-1$
    (i.e. half of the) factors in the leftmost product where $i < j$.
    Then, $f'(x) = \sum_{i=1}^n \frac{f(x)}{x - \alpha_i}$,
    so every addend but the $i$th has a factor of $x-\alpha_i$,
    meaning $f'(\alpha_i) = \prod_{j\ne i}(\alpha_i - \alpha_j)$, so
    \begin{equation*}
	D = (-1)^{n(n-1)/2}\prod_{i=1}^n f'(\alpha_i).
    \end{equation*}

    On the other hand, since $\Gal(L/K)$
    and $\Gal(L/F)$ act transitively on $\{\alpha_1,\dots,\alpha_n\}$,
    \begin{equation*}
	N_{K/F}(f'(\alpha))^{[L:K]} = N_{L/F}(f'(\alpha))
	= \prod_{\sigma\in\Gal(L/F)} \sigma(f'(\alpha))
	= \left(\prod_{i=1}^n f'(\alpha_i)\right)^{[L:K]}
    \end{equation*}
    where the first equality is true because
    $N_{K/F}(f'(\alpha))\in K$,
    so $D = (-1)^{n(n-1)/2} N_{K/F}(f'(\alpha))$.
\end{soln}

\begin{problem}
    [\S14.7\#2]
    Let $\zeta$ be a primitive $7$th root of unity
    and let $\alpha = \zeta + \zeta\inv$.
    \begin{enumerate}[(a)]
	\ii Show that $\zeta$ is a root of
	$z^2 - \alpha z + 1\in \QQ(\alpha)[z]$.
	\ii Show using the minimal polynomial for $\zeta$ that
	$\alpha$ is a root of the cubic
	$x^3 + x^2 - 2x - 1$.
	\ii Use (a) and (b) together with
	the explicit solution of the cubic in (b)
	to express $\zeta$ in terms of radicals.
    \end{enumerate}
\end{problem}

\begin{soln}
    (a) Plugging in $z = \alpha$, we compute
    \begin{equation*}
        \zeta^2 - (\zeta + \zeta\inv)\zeta + 1
	= \zeta^2 - \zeta^2 - 1 + 1 = 0.
    \end{equation*}

    (b) Since $\zeta^6 + \zeta^5 + \zeta^4 + \zeta^3
    + \zeta^2 + \zeta + 1 = 0$, dividing by $\zeta^3$ gives
    \begin{equation*}
	\alpha^3 - \alpha^2 - 2\alpha - 1
	= (\alpha^3 - 3\alpha)
	+ (\alpha^2 - 2)
	+ \alpha + 1
	= (\zeta^3 + \zeta^{-3}) + (\zeta^2 + \zeta^{-2}) +
	(\zeta + \zeta\inv) + 1 = 0.
    \end{equation*}

    (c) no thanks
    (The idea is of course to pick a root of the cubic to be $\alpha$
    and use the fact that $\alpha$ is twice the real part of $\zeta$
    and $\abs\zeta = 1$ to compute the imaginary part of $\zeta$,
    but I am not writing that out.)
\end{soln}

\pagebreak
\begin{problem}
    [\S14.7\#7]
    Let $F$ be a field of characteristic not dividing $n$
    containing the $n$th roots of unity
    and let $K$ be a cyclic extension of degree $d$ dividing $n$.
    Then $K = F[\sqrt[n] a]$ for some nonzero $a\in F$.
    Let $\sigma$ be a generator for $\Gal(K/F)$.
    \begin{enumerate}[(a)]
	\ii Show that $\sigma(\sqrt[n]a) = \zeta\sqrt[n]a$
	for some primitive $d$th root of unity $\zeta$.
	\ii Suppose $K = F[\sqrt[n]a] = F[\sqrt[n]b]$.
	Use (a) to show that
	\[
	    \frac{\sigma(\sqrt[n]a)}{\sqrt[n]a}
	    = \left(\frac{\sigma(\sqrt[n]b)}{\sqrt[n]b}\right)^i
	\]
	for some integer $i$ relatively prime to $d$.
	Conclude that $\sigma$ fixes
	$\frac{\sqrt[n]a}{(\sqrt[n]b)^i}$,
	so this is an element of $F$.
	\ii Prove that $K = F[\sqrt[n]a] = F[\sqrt[n]b]$
	if and only if $a = b^ic^n$
	and $b = a^jd^n$ for some $c,d\in F$,
	i.e. $a$ and $b$ represent the same subgroup
	of $F^\times/(F^\times)^n$.
    \end{enumerate}
\end{problem}

\begin{soln}
    (a) The minimal polynomial of $\sqrt[n]a$ divides $x^n - a$,
    so the Galois conjugates of $\sqrt[n]a$
    are $\sqrt[n]a$ times some $n$th root of unity,
    so $\sigma(\sqrt[n]a) = \zeta\sqrt[n]a$
    for an $n$th root of unity $\zeta$.
    The order of $\zeta$ in $F^\times$
    is equal to the order of $\sigma$ in $\Gal(K/F)$
    because $\sigma$ is determined by its action on $\sqrt[n]a$.
    This order is equal to $d$,
    so $\zeta$ is a primitive $d$th root of unity.
    \\

    (b) By (a), $\sigma(\sqrt[n]a) = \zeta_1\sqrt[n]a$
    and $\sigma(\sqrt[n]b) = \zeta_2\sqrt[n]b$
    for primitive $d$th roots of unity $\zeta_1$ and $\zeta_2$.
    These roots of unity both generate the cyclic group
    of roots of unity $\mu_d$, so some power $\zeta_2^i$ of $\zeta_2$
    must equal $\zeta_1$ as desired.
    Rearrangning the resulting equality then shows
    \begin{equation*}
	\sigma\left(\frac{\sqrt[n]a}{(\sqrt[n]b)^i}\right)
	= \frac{\sigma(\sqrt[n]a)}{(\sigma(\sqrt[n]b))^i}
	= \frac{\sqrt[n]a}{(\sqrt[n]b)^i},
    \end{equation*}
    which means that $\frac{\sqrt[n]a}{(\sqrt[n]b)^i}$
    is fixed by $\sigma$, which only fixes $F$
    because it generates $\Gal(K/F)$.
    \\

    (c) If $a = b^i c^n$ and $b = a^j d^n$
    then $\sqrt[n]a = c(\sqrt[n]b)^i$ implies
    $\sqrt[n]a\in F[\sqrt[n]b]$
    and similarly $\sqrt[n]b = d(\sqrt[n]a)^j$ implies
    $\sqrt[n]b\in F[\sqrt[n]a]$.
    Conversely, if $\sqrt[n]a$ and $\sqrt[n]b$ generate
    the same extension of $F$, then (b) implies that
    $\sqrt[n]a = c(\sqrt[n]b)^i$ for some $c\in F$,
    and applying (b) again after swapping $a$ and $b$
    shows that $\sqrt[n]b = d(\sqrt[n]a)^j$ for some integer $i$
    and $d\in F$.
\end{soln}









\end{document}
